\documentclass{situationist-preprint}

\title{Self-Knowledge of Recalcitrant Attitudes}
\msid{EPI-15-b}
\msfamily{EPI-15}
\msdiscipline{Epistemology}
\mstype{Companion manuscript}
\mskeywords{self-knowledge; transparency; recalcitrant attitudes; alief; self-perception; introspection}

\begin{document}
\maketitle

\begin{abstract}
EPI-15-a shows that the transparent method yields knowledge of answerable attitudes and fails for attitudes that are not answerable to reasons. This companion examines the second class in more detail. We distinguish three kinds of recalcitrant attitude: the fear that persists against one's verdict that there is no danger, the desire that persists against one's verdict that its object is not worth having, and the belief held against one's verdict. For each we identify the route by which it can be known, and show that knowledge so acquired is inferential, fallible and without first-person authority, but is knowledge when its basis is safe.
\end{abstract}

\section{Three kinds of recalcitrance}

An attitude is answerable when the subject's verdict on the question it settles fixes whether she has it (Definition~2.2 of EPI-15-a). Theorem~3.3 there shows that, for an attitude that is not answerable, the transparent self-ascription is unsafe, and Corollary~3.4 that knowledge of it, where possible, comes by other routes. The routes differ by kind.

\emph{Recalcitrant fear.} The nervous flyer judges that flying is safe and is afraid. Fear has a phenomenal profile, and introspection is reliable about current phenomenal states though not about their causes (MND-20). She can therefore know by introspection that she feels afraid. That the fear is a fear \emph{of flying} is not given by the feeling; she learns it by noting that it rises at the gate and subsides on landing.

\emph{Recalcitrant desire.} The smoker who judges that the cigarette is not worth having and wants it has no reliable phenomenal mark of the desire, which may be dispositional. He knows it chiefly from his conduct: from what he reaches for and what he plans around. This is the route that self-perception theory takes as the general case \citep{Bem1972}; it is general for recalcitrant attitudes, not for attitudes in general.

\emph{Belief against one's verdict.} The subject who sincerely judges that the members of a group are as competent as anyone, but whose expectations, attention and choices treat them as less competent, has a state that guides action as belief does and does not answer to her verdict. Whether it is a belief or an alief, a state of a more primitive kind \citep{Gendler2008}, it is not answerable, and it is known, when it is, by inference from patterns in conduct, sometimes with the help of testimony from others who observe those patterns.

\begin{proposition}[Character of recalcitrant self-knowledge]\label{p:char}
Knowledge of a recalcitrant attitude is (i) based on evidence about oneself, phenomenal or behavioural, rather than on a verdict on the question the attitude settles; (ii) fallible in the same way as knowledge of another's attitudes; (iii) without first-person authority; and (iv) knowledge when its basis is safe, by Principle~P1 of EPI-15-a.
\end{proposition}

\begin{proof}[Argument]
(i) and (iii) follow from Theorem~3.3 of EPI-15-a: the verdict does not fix the attitude, so a self-ascription based on it is unsafe, and no ascription is authoritative merely by being the subject's own. (ii) The evidence used is of the same kind that others could use, with the addition of phenomenal evidence in the case of fear. (iv) is an instance of the principle.
\end{proof}

\section{An objection}

\begin{objection}
The subject who discovers, as in Moran's example, that she has long felt betrayed by a friend comes to know an attitude that she can then avow; this is first-personal knowledge of a recalcitrant attitude \citep{Moran2001}.
\end{objection}

\begin{reply}
What she first acquires is knowledge from the theoretical stance, by evidence. Avowal becomes possible only if, on reflection, she endorses the resentment as warranted, and then the attitude has become answerable. If she judges it unwarranted and it persists, the avowal is not available, and her knowledge remains of the kind described by Proposition~\ref{p:char}.
\end{reply}

\section{Conclusion}

Recalcitrant attitudes can be known, but as one knows another's mind: by introspection of feeling where there is a feeling, and otherwise by inference from conduct. The limit of transparency is a limit of a method, not of self-knowledge.

\begin{thebibliography}{99}\small
\bibitem[Bem(1972)]{Bem1972} Bem, D.~J. (1972). Self-perception theory. In L. Berkowitz (Ed.), \emph{Advances in Experimental Social Psychology}, Vol.~6 (pp.~1--62). Academic Press.
\bibitem[Gendler(2008)]{Gendler2008} Gendler, T.~S. (2008). Alief and belief. \emph{The Journal of Philosophy}, 105(10), 634--663.
\bibitem[Moran(2001)]{Moran2001} Moran, R. (2001). \emph{Authority and Estrangement: An Essay on Self-Knowledge}. Princeton University Press.
\end{thebibliography}

\end{document}
